%=============================================================================!
% 2.9  SEVERAL BODIES: MOMENTUM AND THE CENTRE OF MASS                        !
%=============================================================================!
\section{Several bodies: momentum and the centre of mass}
\label{sec:ne-bodies}

So far each body has moved under forces from its surroundings. When
several bodies push and pull on one another, as the atoms of a material
do, the third law ties their motions together. This section shows how,
using a new quantity, momentum, and a new point, the centre of mass.

\subsection*{Momentum}

The momentum of a body is its mass times its velocity,
\begin{equation*}
   \vb{p} = m\vb{v} ,
\end{equation*}
a vector in \si{\kilogram\metre\per\second} that points along the
velocity. Since the mass is constant, the rate of change of the momentum
is $\dd\vb{p}/\dd t = m\,\dd\vb{v}/\dd t = m\vb{a}$, and the second law
can be written
\begin{equation*}
   \frac{\dd\vb{p}}{\dd t} = \vb{F} :
\end{equation*}
the net force on a body is the rate at which its momentum changes. A
constant force $\vb{F}$ acting from $t = 0$ to a time $t$ therefore
changes the momentum by $\vb{F}t$, since $\vb{p}$ then grows at the
constant rate $\vb{F}$, as $v$ grows by $at$ at the constant rate $a$ in
\cref{eq:mo-constant}; so the momentum of a body measures how long
a given force must push to bring it from rest to its velocity, or to stop
it again. A trolley of \SI{20}{\kilogram} pushed from rest by
\SI{10}{\newton} for \SI{2}{\second} gains a momentum of $10\times2 =
\SI{20}{\kilogram\metre\per\second}$, and so a velocity of $20/20 =
\SI{1}{\metre\per\second}$. A lorry and a bicycle moving at the same
speed have very different momenta, and the brakes of the lorry must push
far harder, or for far longer, to stop it.

\subsection*{Two carts}

Two carts, of masses $m_1 = \SI{1}{\kilogram}$ and $m_2 =
\SI{3}{\kilogram}$, stand at rest on a level, smooth track with a
compressed spring between them, and are released
(\cref{fig:ne-carts-sketch}). On each cart the weight and the normal
force balance, so only the forces along the track matter, and these come
from the spring. The spring is light, and we treat it as having no mass.
Two steps then fix the forces on the carts. First, the net force on the
spring must be zero, since by the second law any net force on a body of
no mass would give it an unlimited acceleration; so the two carts push on
the two ends of the spring with forces of equal size. Second, by the third
law, the spring pushes back on each cart with a force of that same size:
$-\vb{F}$ on the first cart, backwards, and $+\vb{F}$ on the second,
forwards.

\begin{figure}[htb]
\centering
\begin{tikzpicture}[line cap=round,
   force/.style={-{Stealth[length=5pt]}, line width=1.1pt, accent},
   coil/.style={decorate, decoration={coil, aspect=0.45,
      segment length=3pt, amplitude=3.5pt, pre length=2pt,
      post length=2pt}, line width=0.6pt}]
   \draw[line width=0.6pt] (-0.6,0) -- (7.4,0);
   % cart 1, the lighter, and cart 2, the heavier
   \draw[line width=0.6pt, fill=black!8] (0.6,0.18) rectangle (1.8,0.9);
   \draw[line width=0.6pt, fill=black!8] (3.6,0.18) rectangle (6.0,1.2);
   \foreach \x in {0.85,1.55,4.0,5.6}
      \draw[line width=0.6pt, fill=white] (\x,0.09) circle (0.09);
   \draw[coil] (1.8,0.5) -- (3.6,0.5);
   \node at (1.2,0.54) {\small 1};
   \node at (4.8,0.69) {\small 2};
   % the push of the spring on each cart
   \draw[force] (0.6,0.54) -- (-0.4,0.54) node[above] {\small $-\vb{F}$};
   \draw[force] (6.0,0.69) -- (7.0,0.69) node[above] {\small $+\vb{F}$};
   \node[below] at (1.2,-0.05) {\small \SI{1}{\kilogram}};
   \node[below] at (4.8,-0.05) {\small \SI{3}{\kilogram}};
\end{tikzpicture}
\caption{Two carts at rest with a compressed spring between them. The
spring pushes the two carts with forces of equal size in opposite
directions (teal); the weights and normal forces, which balance on each
cart, are not drawn.}
\label{fig:ne-carts-sketch}
\end{figure}

Applying the second law to each cart and adding the two equations,
\begin{equation*}
   \frac{\dd\vb{p}_1}{\dd t} = -\vb{F}, \qquad
   \frac{\dd\vb{p}_2}{\dd t} = +\vb{F}
   \quad\Longrightarrow\quad
   \frac{\dd}{\dd t}\bigl(\vb{p}_1 + \vb{p}_2\bigr) = -\vb{F} + \vb{F}
   = \vb{0} .
\end{equation*}
The total momentum $\vb{P} = \vb{p}_1 + \vb{p}_2$ therefore does not
change, however the force of the spring varies, and since the carts
started at rest it stays zero. A quantity that keeps the same value as
time passes is said to be conserved, and here the total momentum is
conserved. The carts therefore always have equal and opposite momenta,
$m_1\vb{v}_1 = -m_2\vb{v}_2$, and the lighter cart moves faster by the
ratio of the masses, three times faster here.

\Cref{fig:ne-carts} shows the motion computed for a spring of stiffness
\SI{200}{\newton\per\metre} compressed by \SI{10}{\centi\metre}. The
spring is not fixed to either cart, so once it reaches its natural length
it could only stay in contact by pulling, which a loose spring cannot do,
and it falls away after about \SI{0.096}{\second}, a time found, like
the speeds below, by a computer solution. From then on no force
acts along the track, and by the first law the carts coast apart at
constant velocities, \SI{-1.225}{\metre\per\second} and
\SI{+0.408}{\metre\per\second}, a ratio of exactly $-3$. Their momenta,
$\mp\SI{1.225}{\kilogram\metre\per\second}$, add to zero at every
instant. Chapter~3 shows how energy gives these speeds directly.

\begin{figure}[htb]
\centering
\includegraphics{ch02_newton/carts}
\caption{The carts of \cref{fig:ne-carts-sketch} pushed apart by a
spring of \SI{200}{\newton\per\metre} compressed by
\SI{10}{\centi\metre}; the spring falls away at the dotted line. (a)
Positions of the carts and of their centre of mass (dashed). (b) Momenta
of the carts and the component $P$ of their total momentum along the
track (dashed), which stays zero.}
\label{fig:ne-carts}
\end{figure}

\subsection*{Many bodies}

The same argument works for any number of bodies. The forces on each body
are of two kinds. External forces come from outside the group, such as
the weight of a cart, and we write the net external force on body $i$ as
$\vb{F}_i^{\mathrm{ext}}$. Internal forces come from the other members,
and we write the force on body $i$ from body $j$ as $\vb{F}_{ij}$. For
three bodies, with total momentum $\vb{P} = \vb{p}_1 + \vb{p}_2 +
\vb{p}_3$, adding the three second laws and then regrouping the
internal forces by pairs gives
\begin{align*}
   \frac{\dd\vb{P}}{\dd t}
   &= \frac{\dd\vb{p}_1}{\dd t} + \frac{\dd\vb{p}_2}{\dd t}
      + \frac{\dd\vb{p}_3}{\dd t}\\
   &= \bigl(\vb{F}_1^{\mathrm{ext}} + \vb{F}_{12} + \vb{F}_{13}\bigr)
      + \bigl(\vb{F}_2^{\mathrm{ext}} + \vb{F}_{21} + \vb{F}_{23}\bigr)
      + \bigl(\vb{F}_3^{\mathrm{ext}} + \vb{F}_{31} + \vb{F}_{32}\bigr)\\
   &= \vb{F}^{\mathrm{ext}}
      + \bigl(\vb{F}_{12} + \vb{F}_{21}\bigr)
      + \bigl(\vb{F}_{13} + \vb{F}_{31}\bigr)
      + \bigl(\vb{F}_{23} + \vb{F}_{32}\bigr) ,
\end{align*}
where $\vb{F}^{\mathrm{ext}} = \vb{F}_1^{\mathrm{ext}} +
\vb{F}_2^{\mathrm{ext}} + \vb{F}_3^{\mathrm{ext}}$ is the total external
force. By the third law each bracket is zero. With $N$ bodies, numbered
$i = 1, \dots, N$, the total momentum is $\vb{P} = \vb{p}_1 + \dotsb +
\vb{p}_N$, written $\sum_{i=1}^{N}\vb{p}_i$ with the summation sign of
\cref{sec:mo-integrating}, and the regrouping is the same: every internal force $\vb{F}_{ij}$ meets its
partner $\vb{F}_{ji}$, each pair cancels, and
\begin{equation}
   \frac{\dd\vb{P}}{\dd t} = \vb{F}^{\mathrm{ext}} .
   \label{eq:ne-momentum}
\end{equation}
The internal forces, however strong and however complicated, cannot
change the total momentum. Whenever the total external force is zero,
the total momentum is conserved. That happens when the external forces
cancel, as the weights and normal forces of the carts do, and also when
there are no external forces at all, in which case the group is called
isolated.

\subsection*{The centre of mass}

The group as a whole has a position of its own, the centre of mass, an
average of the positions of its members in which each counts in
proportion to its mass,
\begin{equation*}
   \vb{R} = \frac{1}{M}\sum_{i=1}^{N} m_i\vb{r}_i, \qquad
   M = \sum_{i=1}^{N} m_i ,
\end{equation*}
where $M$ is the total mass. For two bodies on a line it is the balance
point: a light rod carrying the two carts would balance on a pivot placed
there. The centres of the carts start at $x = 0$ and $x =
\SI{0.2}{\metre}$, so their
centre of mass is at $(1\times0 + 3\times0.2)/4 = \SI{0.15}{\metre}$,
nearer the heavier cart. Its motion follows in three lines,
\begin{align*}
   M\vb{R} &= \sum_{i=1}^{N} m_i\vb{r}_i
      && \text{multiplying by $M$}\\
   M\,\frac{\dd\vb{R}}{\dd t} &= \sum_{i=1}^{N} m_i\vb{v}_i = \vb{P}
      && \text{differentiating term by term, the masses constant}\\
   M\,\frac{\dd^2\vb{R}}{\dd t^2} &= \frac{\dd\vb{P}}{\dd t}
      = \vb{F}^{\mathrm{ext}}
      && \text{differentiating again, by \cref{eq:ne-momentum}.}
\end{align*}
The centre of mass moves as a single body of mass $M$ would under the
external forces alone. For the carts there is no external force along the
track, so the centre of mass stays at \SI{0.15}{\metre} while the carts
fly apart (\cref{fig:ne-carts}a). A firework that bursts in flight shows
the same thing in two dimensions: the pieces scatter in every direction,
but until the first of them lands, their centre of mass carries on along
the parabola that the unburst firework would have followed. The only
external force is still the weight, $\vb{F}^{\mathrm{ext}} = \sum_i
m_i\vb{g} = M\vb{g}$, so $\dd^2\vb{R}/\dd t^2 = \vb{g}$ as before the
burst, and the internal forces of the burst change neither $\vb{R}$ nor
$\dd\vb{R}/\dd t = \vb{P}/M$. Chapter~7 shows that the
conservation of momentum follows from a symmetry, that the laws of motion
are the same at every place.

\begin{keyidea}
Internal forces cancel in pairs by the third law, so the total momentum
of a group changes only through external forces, and is conserved when
they add to zero. The centre of mass moves as a single body under the
external forces alone.
\end{keyidea}

\notebook{02}{change the mass ratio of the carts; the centre of mass never
moves}

\begin{selfcheck}
A skater of \SI{60}{\kilogram} standing on smooth ice throws a ball of
\SI{0.5}{\kilogram} forwards at \SI{12}{\metre\per\second}, measured
relative to the ice. What is the
skater's velocity afterwards?
\end{selfcheck}
